(1) 可逆绝热满足 \(TV^{\gamma-1}=\mathrm{const}\)，所以 \(\boxed{T_2=T_1r^{\gamma-1}}\)；膨胀从小体积到大体积给 \(\boxed{T_4=T_3/r^{\gamma-1}}\)。所有压强随后可由 \(pV=RT\) 求出，闭合条件是同一初末状态。

(2) 定容加热的输入热量为 \(Q_h=C_V(T_3-T_2)>0\)，定容排热量的正值是 \(Q_c=C_V(T_4-T_1)>0\)。一循环内内能变化为零，故净功 \(W=Q_h-Q_c\)，热效率 \(\eta=W/Q_h=1-(T_4-T_1)/(T_3-T_2)\)。

(3) 代入绝热温度关系得 \(T_4-T_1=(T_3-T_2)/r^{\gamma-1}\)，所以 \(\boxed{\eta=1-r^{1-\gamma}}\)。固定 \(\gamma>1\) 时对 \(r\) 求导 \(d\eta/dr=(\gamma-1)r^{-\gamma}>0\)，形式上 \(r\to\infty\) 效率趋1；但固定峰温时大压缩比会让加热条件 \(T_3>T_1r^{\gamma-1}\) 失效。

(4) 理想循环中两绝热过程熵守恒，两定容换热过程的气体熵变分别为 \(C_V\ln(T_3/T_2)\)、\(C_V\ln(T_1/T_4)\)，两者和为零。真实机器有限时间燃烧、摩擦、排气、热漏与压缩末端自燃爆震限制压缩比，不能将数学效率极限当成实际可实现状态。
